% Source-faithful mathematical blueprint of Hartshorne, Algebraic Geometry (1977), Chapter I.
% Authorial exposition, examples, and exercises are omitted; definitions, statements, and proof
% arguments are retained. Source coordinates are the printed Chapter I coordinates.

\section{Affine Varieties}

Fix an algebraically closed field $k$, put $A=k[x_1,\ldots,x_n]$, and write
\(\mathbb A^n=k^n\).  For $T\subset A$, set
\[
Z(T)=\{P\in\mathbb A^n:f(P)=0\text{ for every }f\in T\}.
\]

\begin{definition}[Polynomial evaluation and zero sets]
\label{def:ha-ch1-polynomial-evaluation}
\source{hartshorne-algebraic-geometry:I.1}
\lean{Hartshorne.AffinePoint, Hartshorne.AffinePolynomial, Hartshorne.evaluate,
Hartshorne.zeroSet, Hartshorne.commonZeroSet, Hartshorne.commonZeroSet_span,
Hartshorne.commonZeroSet_span_finite}
\leanok
For $P=(a_1,\ldots,a_n)\in\mathbb A^n$ and $f\in A$, regard $f$ as the function
$f(P)=f(a_1,\ldots,a_n)$.  For one polynomial write
$Z(f)=\{P\in\mathbb A^n:f(P)=0\}$, and for a subset $T\subseteq A$ write $Z(T)$
for the common zero set displayed above.  If $\mathfrak a=(T)$, then
$Z(T)=Z(\mathfrak a)$; since $A$ is noetherian, $\mathfrak a$ has finitely many generators.
\end{definition}

\begin{definition}[Algebraic set]
\label{def:ha-ch1-algebraic-set}
\source{hartshorne-algebraic-geometry:I.1}
\lean{Hartshorne.IsAlgebraicSet, Hartshorne.isAlgebraicSet_iff}
\leanok
A subset $Y\subseteq\mathbb A^n$ is algebraic if and only if $Y=Z(T)$ for some subset
$T\subseteq A$.
\uses{def:ha-ch1-polynomial-evaluation}
\end{definition}

\begin{proposition}[Affine algebraic sets]
\label{prop:ha-ch1-1-1}
\source{hartshorne-algebraic-geometry:I.1.1}
\dcref{I.1.1}
The union of two algebraic sets and every intersection of algebraic sets are algebraic; so are
$\varnothing$ and $\mathbb A^n$.
\uses{def:ha-ch1-algebraic-set,def:ha-ch1-polynomial-evaluation}
\end{proposition}
\begin{proof}
If $Y_i=Z(T_i)$, then $Y_1\cup Y_2=Z(T_1T_2)$, where $T_1T_2$ consists of all products
of one element from each set.  Also $\bigcap_\alpha Z(T_\alpha)=Z(\bigcup_\alpha T_\alpha)$,
while $Z(1)=\varnothing$ and $Z(0)=\mathbb A^n$.
\end{proof}

\begin{definition}[Zariski topology]
\label{def:ha-ch1-zariski-topology}
\source{hartshorne-algebraic-geometry:I.1}
The Zariski topology on $\mathbb A^n$ is the topology whose closed subsets are the algebraic
sets; equivalently, the complements of the sets $Z(T)$ are the open subsets.
\uses{def:ha-ch1-algebraic-set}
\end{definition}

\begin{definition}[Irreducible subset]
\label{def:ha-ch1-irreducible-subset}
\source{hartshorne-algebraic-geometry:I.1}
A nonempty subset $Y$ of a topological space is irreducible if it is not the union of two proper
closed subsets (with the induced topology). The empty set is not irreducible.
\end{definition}

\begin{definition}[Affine algebraic variety]
\label{def:ha-ch1-affine-variety}
\source{hartshorne-algebraic-geometry:I.1}
An affine algebraic variety is an irreducible algebraic set.
\uses{def:ha-ch1-algebraic-set,def:ha-ch1-irreducible-subset}
\end{definition}

\begin{definition}[Quasi-affine variety]
\label{def:ha-ch1-quasi-affine-variety}
\source{hartshorne-algebraic-geometry:I.1}
A quasi-affine variety is an open subset of an affine algebraic variety.
\uses{def:ha-ch1-affine-variety}
\end{definition}

\begin{definition}[Vanishing ideal]
\label{def:ha-ch1-vanishing-ideal}
\source{hartshorne-algebraic-geometry:I.1}
\lean{Hartshorne.vanishingIdeal, Hartshorne.mem_vanishingIdeal}
\leanok
For $Y\subseteq\mathbb A^n$, its ideal is
\[
 I(Y)=\{f\in A:f(P)=0\text{ for every }P\in Y\}.
\]
\uses{def:ha-ch1-polynomial-evaluation}
\end{definition}

\begin{proposition}[Ideal-zero-set identities]
\label{prop:ha-ch1-1-2}
\source{hartshorne-algebraic-geometry:I.1.2}
\dcref{I.1.2}
For $T_1\subseteq T_2\subset A$, $Z(T_1)\supseteq Z(T_2)$; for $Y_1\subseteq Y_2$,
$I(Y_1)\supseteq I(Y_2)$; $I(Y_1\cup Y_2)=I(Y_1)\cap I(Y_2)$; for an ideal $\mathfrak a$,
$I(Z(\mathfrak a))=\sqrt{\mathfrak a}$; and $Z(I(Y))=\overline Y$.
\uses{def:ha-ch1-polynomial-evaluation,def:ha-ch1-algebraic-set,def:ha-ch1-vanishing-ideal,thm:ha-ch1-1-3A}
\end{proposition}
\begin{proof}
The first three assertions follow immediately from the definitions.  The fourth is the
Nullstellensatz (Theorem~\ref{thm:ha-ch1-1-3A}).  For the last, $Y\subseteq Z(I(Y))$; if
$W=Z(\mathfrak a)$ is any closed set containing $Y$, then $\mathfrak a\subseteq I(W)\subseteq I(Y)$,
so $W\supseteq Z(I(Y))$.  Thus $Z(I(Y))$ is the smallest closed set containing $Y$.
\end{proof}

\begin{theorem}[Hilbert Nullstellensatz]
\label{thm:ha-ch1-1-3A}
\source{hartshorne-algebraic-geometry:I.1.3A}
\dcref{I.1.3A}
\lean{Hartshorne.hilbertNullstellensatz}
\leanok
If $\mathfrak a\subset A$ and $f$ vanishes on $Z(\mathfrak a)$, then $f^r\in\mathfrak a$
for some integer $r>0$.
\end{theorem}
\begin{proof}
Hartshorne gives external references for this theorem and does not prove it in Chapter I.
\end{proof}

\begin{corollary}[Affine Nullstellensatz correspondence]
\label{cor:ha-ch1-1-4}
\source{hartshorne-algebraic-geometry:I.1.4}
\dcref{I.1.4}
\lean{Hartshorne.AffineAlgebraicSet, Hartshorne.AffineRadicalIdeal,
Hartshorne.affineAlgebraicSetToRadicalIdeal,
Hartshorne.affineRadicalIdealToAlgebraicSet,
Hartshorne.affineRadicalIdealToAlgebraicSet_leftInverse,
Hartshorne.affineAlgebraicSetToRadicalIdeal_rightInverse,
Hartshorne.affineNullstellensatzEquiv,
Hartshorne.affineAlgebraicSetToRadicalIdeal_antitone,
Hartshorne.affineRadicalIdealToAlgebraicSet_antitone,
Hartshorne.isAffineVariety_iff_vanishingIdeal_isPrime}
\leanok
Algebraic sets correspond inclusion-reversingly to radical ideals by $Y\mapsto I(Y)$ and
$\mathfrak a\mapsto Z(\mathfrak a)$.  An algebraic set is irreducible exactly when its ideal is prime.
\uses{thm:ha-ch1-1-3A,prop:ha-ch1-1-2}
\end{corollary}
\begin{proof}
The first assertion is Theorem~\ref{thm:ha-ch1-1-3A} together with Proposition~\ref{prop:ha-ch1-1-2}.
If $Y$ is irreducible and $fg\in I(Y)$, then $Y\subseteq Z(f)\cup Z(g)$, so one of $f,g$
vanishes on $Y$; hence $I(Y)$ is prime.  Conversely, if $\mathfrak p$ is prime and
$Z(\mathfrak p)=Y_1\cup Y_2$, then $\mathfrak p=I(Y_1)\cap I(Y_2)$, so primality forces
$\mathfrak p=I(Y_i)$ for some $i$, and $Z(\mathfrak p)=Y_i$.
\end{proof}

\begin{definition}[Affine coordinate ring]
\label{def:ha-ch1-affine-coordinate}
\source{hartshorne-algebraic-geometry:I.1.4}
For an affine algebraic set $Y$, put $A(Y)=A/I(Y)$.
\uses{def:ha-ch1-algebraic-set,def:ha-ch1-vanishing-ideal}
\end{definition}

\begin{proposition}[Irreducible components]
\label{prop:ha-ch1-1-5}
\source{hartshorne-algebraic-geometry:I.1.5}
\dcref{I.1.5}
In a noetherian topological space, every nonempty closed set is a finite union of irreducible
closed subsets, uniquely after omitting components contained in another.
\uses{def:ha-ch1-zariski-topology,def:ha-ch1-irreducible-subset}
\end{proposition}
\begin{proof}
If a minimal counterexample $Y$ were reducible, write $Y=Y'\cup Y''$ with both proper closed;
minimality decomposes each into finitely many irreducible closed sets, contradiction.  For
uniqueness, an irreducible component in one decomposition is contained in a component of the
other; applying this in both directions and removing the first common component gives induction.
\end{proof}

\begin{corollary}[Affine components]
\label{cor:ha-ch1-1-6}
\source{hartshorne-algebraic-geometry:I.1.6}
\dcref{I.1.6}
Every affine algebraic set is a unique irredundant union of affine varieties.
\uses{prop:ha-ch1-1-5,cor:ha-ch1-1-4}
\end{corollary}
\begin{proof}
Apply Proposition~\ref{prop:ha-ch1-1-5} to the noetherian Zariski space $Y$.  Its irreducible
components are closed in $Y$, hence affine algebraic sets, and their ideals are prime by
Corollary~\ref{cor:ha-ch1-1-4}.  The uniqueness and irredundancy are exactly the uniqueness clause
in Proposition~\ref{prop:ha-ch1-1-5}.
\end{proof}

For a topological space, dimension is the supremum of lengths of chains of distinct irreducible
closed subsets.  For a ring, the height of a prime is the supremum of lengths of chains ending in
that prime, and the Krull dimension is the supremum of prime heights.

\begin{proposition}[Dimension of an affine coordinate ring]
\label{prop:ha-ch1-1-7}
\source{hartshorne-algebraic-geometry:I.1.7}
\dcref{I.1.7}
For an affine algebraic set $Y$, $\dim Y=\dim A(Y)$.
\uses{def:ha-ch1-affine-coordinate}
\end{proposition}
\begin{proof}
Closed irreducible subsets of $Y$ correspond to prime ideals of $A$ containing $I(Y)$, hence
to prime ideals of $A(Y)$; chains correspond and have the same lengths.
\end{proof}

\begin{theorem}[Dimension theorem for finite type domains]
\label{thm:ha-ch1-1-8A}
\source{hartshorne-algebraic-geometry:I.1.8A}
\dcref{I.1.8A}
If $B$ is a finitely generated domain over a field $k$, then
$\dim B=\operatorname{trdeg}_kK(B)$ and
$\operatorname{ht}\mathfrak p+\dim(B/\mathfrak p)=\dim B$ for every prime $\mathfrak p$.
\end{theorem}
\begin{proof}
Hartshorne cites Matsumura and Atiyah--Macdonald for this commutative-algebra theorem.
\end{proof}

\begin{proposition}[Dimension of affine space]
\label{prop:ha-ch1-1-9}
\source{hartshorne-algebraic-geometry:I.1.9}
\dcref{I.1.9}
\(\dim\mathbb A^n=n\).
\uses{prop:ha-ch1-1-7,thm:ha-ch1-1-8A}
\end{proposition}
\begin{proof}
Apply Proposition~\ref{prop:ha-ch1-1-7} and Theorem~\ref{thm:ha-ch1-1-8A} to the polynomial ring.
\end{proof}

\begin{proposition}[Dimension of a quasi-affine variety]
\label{prop:ha-ch1-1-10}
\source{hartshorne-algebraic-geometry:I.1.10}
\dcref{I.1.10}
If $Y$ is quasi-affine, then $\dim Y=\dim\overline Y$.
\uses{thm:ha-ch1-1-8A}
\end{proposition}
\begin{proof}
Closing a chain gives $\dim Y\le\dim\overline Y$.  Choose a maximal chain in $Y$ beginning at a
point $P$.  Its closures form a maximal chain in $\overline Y$, and the corresponding prime
ideals lie under the maximal ideal $\mathfrak m_P$ with residue field $k$.  Theorem~\ref{thm:ha-ch1-1-8A}(b)
then gives the same chain length as $\dim A(\overline Y)$, which is $\dim\overline Y$.
\end{proof}

\begin{theorem}[Krull Hauptidealsatz]
\label{thm:ha-ch1-1-11A}
\source{hartshorne-algebraic-geometry:I.1.11A}
\dcref{I.1.11A}
If $A$ is noetherian and $f$ is neither a zero divisor nor a unit, every minimal prime containing
$f$ has height one.
\end{theorem}
\begin{proof}
Hartshorne cites Atiyah--Macdonald for the proof.
\end{proof}

\begin{proposition}[Height-one criterion for a UFD]
\label{prop:ha-ch1-1-12A}
\source{hartshorne-algebraic-geometry:I.1.12A}
\dcref{I.1.12A}
A noetherian domain is a unique factorization domain iff every height-one prime is principal.
\end{proposition}
\begin{proof}
Hartshorne cites Matsumura and Bourbaki for the proof.
\end{proof}

\begin{proposition}[Affine hypersurfaces]
\label{prop:ha-ch1-1-13}
\source{hartshorne-algebraic-geometry:I.1.13}
\dcref{I.1.13}
A variety $Y\subset\mathbb A^n$ has dimension $n-1$ iff $Y=Z(f)$ for one nonconstant irreducible
$f\in k[x_1,\ldots,x_n]$.
\uses{thm:ha-ch1-1-11A,prop:ha-ch1-1-12A}
\end{proposition}
\begin{proof}
For irreducible $f$, $(f)$ is prime and has height one by Hauptidealsatz, so the dimension theorem
gives $n-1$.  Conversely, $Y$ corresponds to a height-one prime in the polynomial UFD; by the
height-one criterion it is generated by an irreducible $f$, hence $Y=Z(f)$.
\end{proof}

\section{Projective Varieties}

Let $S=k[x_0,\ldots,x_n]$ with its total-degree grading.

\begin{definition}[Projective space and homogeneous coordinates]
\label{def:ha-ch1-projective-space}
\source{hartshorne-algebraic-geometry:I.2}
The projective space $\mathbb P^n$ is the set of nonzero $(n+1)$-tuples in $k^{n+1}$ modulo
scalar multiplication. A homogeneous polynomial has all terms of one total degree, and $S$ is
graded by total degree.
\uses{def:ha-ch1-polynomial-evaluation}
\end{definition}

\begin{definition}[Projective algebraic set]
\label{def:ha-ch1-projective-algebraic-set}
\source{hartshorne-algebraic-geometry:I.2}
For homogeneous $T\subset S$, define $Z(T)\subset\mathbb P^n$ by common vanishing. A projective
algebraic set is a set $Z(T)$; its homogeneous vanishing ideal is $I(Y)$ and its homogeneous
coordinate ring is $S(Y)=S/I(Y)$.
\uses{def:ha-ch1-projective-space}
\end{definition}

\begin{definition}[Projective and quasi-projective varieties]
\label{def:ha-ch1-projective-variety}
\source{hartshorne-algebraic-geometry:I.2}
A projective variety is an irreducible projective algebraic set. A quasi-projective variety is an
open subset of a projective variety.
\uses{def:ha-ch1-projective-algebraic-set,def:ha-ch1-irreducible-subset}
\end{definition}

\begin{proposition}[Projective algebraic sets]
\label{prop:ha-ch1-2-1}
\source{hartshorne-algebraic-geometry:I.2.1}
\dcref{I.2.1}
Unions and arbitrary intersections of projective algebraic sets are algebraic, as are $\varnothing$
and $\mathbb P^n$.
\uses{prop:ha-ch1-1-1}
\end{proposition}
\begin{proof}
Hartshorne leaves the proof to the reader; it is the homogeneous analogue of Proposition 1.1.
\end{proof}

\begin{proposition}[Standard projective charts]
\label{prop:ha-ch1-2-2}
\source{hartshorne-algebraic-geometry:I.2.2}
\dcref{I.2.2}
With $H_i=Z(x_i)$ and $U_i=\mathbb P^n-H_i$, the ratio map
$\varphi_i:U_i\to\mathbb A^n$ is a homeomorphism for the induced/projective and affine Zariski
topologies.
\uses{def:ha-ch1-projective-space,def:ha-ch1-projective-algebraic-set,def:ha-ch1-zariski-topology}
\end{proposition}
\begin{proof}
For $i=0$, dehomogenization sends a homogeneous polynomial $f$ to $f(1,y_1,\ldots,y_n)$;
homogenization sends a polynomial $g$ of degree $d$ to
$x_0^d g(x_1/x_0,\ldots,x_n/x_0)$.  These maps identify projective closed subsets in $U_0$
with affine closed subsets and are inverse, proving the claim.  Other charts are identical.
\end{proof}

\begin{corollary}[Affine chart cover]
\label{cor:ha-ch1-2-3}
\source{hartshorne-algebraic-geometry:I.2.3}
\dcref{I.2.3}
Every projective (respectively quasi-projective) variety is covered by its intersections with the
$U_i$, which are affine (respectively quasi-affine) varieties.
\uses{prop:ha-ch1-2-2}
\end{corollary}
\begin{proof}
The $U_i$ cover projective space because a nonzero homogeneous coordinate tuple has some nonzero
coordinate.  Proposition~\ref{prop:ha-ch1-2-2} identifies $Y\cap U_i$ with a closed subset of
$\mathbb A^n$ when $Y$ is projective, and with an open subset of such a closed set when $Y$ is
quasi-projective.
\end{proof}

\section{Morphisms}

\begin{definition}[Regular function]
\label{def:ha-ch1-regular-function}
\source{hartshorne-algebraic-geometry:I.3}
A function to $k$ is regular at $P$ if on a neighborhood of $P$ it is $g/h$ with $h$ nowhere
zero; in projective space $g,h$ are homogeneous of the same degree. It is regular if it is regular
at every point.
\uses{def:ha-ch1-quasi-affine-variety,def:ha-ch1-projective-variety}
\end{definition}

\begin{definition}[Morphism and isomorphism]
\label{def:ha-ch1-morphism}
\source{hartshorne-algebraic-geometry:I.3}
A morphism is a continuous map pulling regular functions back to regular functions. An isomorphism
is a morphism having a morphism as inverse.
\uses{def:ha-ch1-regular-function}
\end{definition}

\begin{lemma}[Continuity of regular functions]
\label{lem:ha-ch1-3-1}
\source{hartshorne-algebraic-geometry:I.3.1}
\dcref{I.3.1}
Every regular function is continuous.
\end{lemma}
\begin{proof}
It suffices to show that each fiber over $a\in k$ is closed.  On a neighborhood where $f=g/h$,
with $h$ nowhere zero, the fiber is the zero set of $g-ah$; local closedness therefore implies
global closedness.
\end{proof}

\begin{theorem}[Affine variety functions and local rings]
\label{thm:ha-ch1-3-2}
\source{hartshorne-algebraic-geometry:I.3.2}
\dcref{I.3.2}
For affine $Y$, $\mathcal O(Y)\cong A(Y)$; points correspond to maximal ideals; for $P$,
$\mathcal O_{P,Y}\cong A(Y)_{\mathfrak m_P}$ and $\dim\mathcal O_{P,Y}=\dim Y$; and
$K(Y)$ is the fraction field of $A(Y)$ with transcendence degree $\dim Y$.
\uses{def:ha-ch1-affine-coordinate,def:ha-ch1-affine-variety,cor:ha-ch1-1-4,thm:ha-ch1-1-8A}
\end{theorem}
\begin{proof}
Polynomial restriction gives an injective map $A(Y)\to\mathcal O(Y)$.  The affine Nullstellensatz
identifies points with maximal ideals and gives the stated $\mathfrak m_P$.  Local expressions $g/h$
give the surjection $A(Y)_{\mathfrak m_P}\to\mathcal O_{P,Y}$, while injectivity follows from
the preceding injection.  The dimension and transcendence assertions follow from the dimension
theorem.  Finally, inside the fraction field, a domain is the intersection of its localizations at
all maximal ideals; since $\mathcal O(Y)$ lies between $A(Y)$ and this intersection, equality follows.
\end{proof}

\begin{proposition}[Projective chart isomorphism]
\label{prop:ha-ch1-3-3}
\source{hartshorne-algebraic-geometry:I.3.3}
\dcref{I.3.3}
The chart map $\varphi_i:U_i\to\mathbb A^n$ of Proposition 2.2 is an isomorphism of varieties.
\uses{prop:ha-ch1-2-2}
\end{proposition}
\begin{proof}
The dehomogenization and homogenization maps identify local quotients of equal-degree homogeneous
polynomials with local quotients of affine polynomials, so the homeomorphism preserves regular
functions in both directions.
\end{proof}

\begin{theorem}[Functions on a projective variety]
\label{thm:ha-ch1-3-4}
\source{hartshorne-algebraic-geometry:I.3.4}
\dcref{I.3.4}
For projective $Y$ with homogeneous coordinate ring $S(Y)$, $\mathcal O(Y)=k$,
$\mathcal O_{P,Y}=S(Y)_{(\mathfrak m_P)}$, and $K(Y)$ is the degree-zero fraction field
$S(Y)_{((0))}$.
\uses{def:ha-ch1-projective-variety,def:ha-ch1-projective-algebraic-set}
\end{theorem}
\begin{proof}
Chart localization gives the local-ring and function-field assertions.  If a global function $f$
has chart expressions, choose one common degree $N$ so that multiplication by $x_i^N$ sends it
into $S(Y)$.  Consequently $S(Y)[f]$ is a finite $S(Y)$-module, so $f$ is integral over $S(Y)$.
Taking degree-zero parts makes $f$ integral over $S(Y)_0=k$; algebraic closedness gives $f\in k$.
\end{proof}

\begin{proposition}[Maps to affine varieties]
\label{prop:ha-ch1-3-5}
\source{hartshorne-algebraic-geometry:I.3.5}
\dcref{I.3.5}
For any variety $X$ and affine variety $Y$, morphisms $X\to Y$ correspond naturally to
$k$-algebra homomorphisms $A(Y)\to\mathcal O(X)$.
\uses{lem:ha-ch1-3-6}
\end{proposition}
\begin{proof}
Pullback gives the homomorphism.  Conversely, choose affine coordinates for $Y$ and take the
images of their coordinate functions; they define a map to affine space, and the ideal of $Y$
vanishes after applying the homomorphism, so the image lies in $Y$.  Lemma 3.6 makes the map a
morphism, and the two constructions are inverse.
\end{proof}

\begin{lemma}[Coordinate criterion for affine morphisms]
\label{lem:ha-ch1-3-6}
\source{hartshorne-algebraic-geometry:I.3.6}
\dcref{I.3.6}
A map $\psi:X\to Y\subset\mathbb A^n$ is a morphism iff every coordinate function $x_i\circ\psi$
is regular on $X$.
\end{lemma}
\begin{proof}
Necessity is the definition.  For sufficiency, polynomial pullbacks are regular and therefore
continuous; local quotients of polynomials pull back to regular functions, so $\psi$ is a morphism.
\end{proof}

\begin{corollary}[Affine anti-equivalence]
\label{cor:ha-ch1-3-7}
\source{hartshorne-algebraic-geometry:I.3.7}
\dcref{I.3.7}
Affine varieties are isomorphic iff their coordinate rings are isomorphic as $k$-algebras.
\uses{thm:ha-ch1-3-2,prop:ha-ch1-3-5,lem:ha-ch1-3-6}
\end{corollary}
\begin{proof}
A morphism induces pullback on global regular functions, and Theorem~\ref{thm:ha-ch1-3-2} identifies
these rings with the affine coordinate rings.  Conversely an isomorphism of coordinate rings gives,
by Proposition~\ref{prop:ha-ch1-3-5}, morphisms in both directions.  The two composites pull back to
the identity on coordinate functions, hence are the identity by Lemma~\ref{lem:ha-ch1-3-6}.
\end{proof}

\begin{corollary}[Affine varieties and finite type domains]
\label{cor:ha-ch1-3-8}
\source{hartshorne-algebraic-geometry:I.3.8}
\dcref{I.3.8}
$X\mapsto A(X)$ is an arrow-reversing equivalence between affine varieties over $k$ and finitely
generated integral $k$-domains.
\uses{cor:ha-ch1-1-4,prop:ha-ch1-3-5}
\end{corollary}
\begin{proof}
For an affine variety, $A(X)$ is finitely generated and a domain by the affine Nullstellensatz.
Conversely, write a finitely generated domain as $k[x_1,\ldots,x_n]/\mathfrak p$ with $\mathfrak p$
prime; Corollary~\ref{cor:ha-ch1-1-4} makes $Z(\mathfrak p)$ an affine variety with that coordinate
ring.  Proposition~\ref{prop:ha-ch1-3-5} identifies the arrows contravariantly.
\end{proof}

\begin{theorem}[Finiteness of integral closure]
\label{thm:ha-ch1-3-9A}
\source{hartshorne-algebraic-geometry:I.3.9A}
\dcref{I.3.9A}
If $A$ is a finitely generated domain over $k$, $L$ is finite algebraic over $K(A)$, and $A'$ is
the integral closure of $A$ in $L$, then $A'$ is a finite $A$-module and a finitely generated
$k$-algebra.
\end{theorem}
\begin{proof}
Hartshorne cites Zariski--Samuel and gives no proof in Chapter I.
\end{proof}

\section{Rational Maps}

\begin{definition}[Rational and birational maps]
\label{def:ha-ch1-rational-map}
\source{hartshorne-algebraic-geometry:I.4}
A rational map $X\dashrightarrow Y$ is an equivalence class of morphisms from nonempty open
subsets of $X$, where two representatives agree after restriction to a common nonempty open.
It is dominant if its image is dense, and birational if it has an inverse rational map.
\end{definition}

\begin{definition}[Function field]
\label{def:ha-ch1-function-field}
\source{hartshorne-algebraic-geometry:I.4}
For an irreducible variety $X$, its function field $K(X)$ is the field of rational functions,
equivalently the fraction field of the coordinate ring of any dense affine open subset.
\end{definition}

\begin{lemma}[Density determines morphisms]
\label{lem:ha-ch1-4-1}
\source{hartshorne-algebraic-geometry:I.4.1}
\dcref{I.4.1}
Two morphisms from $X$ to $Y$ agreeing on a nonempty open subset agree everywhere.
\end{lemma}
\begin{proof}
Embed $Y$ in projective space.  The product map to $\mathbb P^n\times\mathbb P^n$ has image of
the dense open in the diagonal, which is closed (the equations $x_i y_j=x_j y_i$).  Hence its
whole image lies in the diagonal.
\end{proof}

\begin{lemma}[Principal open as affine]
\label{lem:ha-ch1-4-2}
\source{hartshorne-algebraic-geometry:I.4.2}
\dcref{I.4.2}
If $Y=V(f)\subset\mathbb A^n$, then $\mathbb A^n-Y$ is isomorphic to
$H=V(x_{n+1}f-1)\subset\mathbb A^{n+1}$, and its coordinate ring is $k[x_1,\ldots,x_n]_f$.
\uses{lem:ha-ch1-3-6}
\end{lemma}
\begin{proof}
Projection from $H$ has inverse $(a_1,\ldots,a_n)\mapsto(a_1,\ldots,a_n,1/f(a))$; Lemma 3.6
shows both maps are morphisms.
\end{proof}

\begin{proposition}[Affine basis]
\label{prop:ha-ch1-4-3}
\source{hartshorne-algebraic-geometry:I.4.3}
\dcref{I.4.3}
Every variety has a basis of open affine subsets.
\uses{lem:ha-ch1-4-2}
\end{proposition}
\begin{proof}
Reduce to quasi-affine $Y\subset\mathbb A^n$.  For $P\in U\subset Y$, choose $f\in I(Ybar-Y)$
with $f(P)\ne0$.  Then $Y-f^{-1}(0)$ is an affine closed subset of the affine principal open
$\mathbb A^n-f^{-1}(0)$ by Lemma 4.2, and lies in $U$ after choosing $f$ from the complement.
\end{proof}

\begin{theorem}[Rational maps and function fields]
\label{thm:ha-ch1-4-4}
\source{hartshorne-algebraic-geometry:I.4.4}
\dcref{I.4.4}
Dominant rational maps $X\dashrightarrow Y$ correspond bijectively to $k$-embeddings
$K(Y)\hookrightarrow K(X)$, giving an arrow-reversing equivalence with finitely generated field
extensions of $k$.
\uses{prop:ha-ch1-3-5,lem:ha-ch1-4-1}
\end{theorem}
\begin{proof}
Pullback on a dense domain gives the embedding.  Conversely, choose affine generators of $K(Y)$;
after shrinking $X$ they are regular, and Proposition 3.5 gives a morphism to an affine open of
$Y$.  Lemma 4.1 makes the resulting rational map independent of choices.  A finitely generated
field is the function field of the affine variety defined by the relations among a finite set of
field generators, proving essential surjectivity and compatibility with composition.
\end{proof}

\begin{corollary}[Birational tests]
\label{cor:ha-ch1-4-5}
\source{hartshorne-algebraic-geometry:I.4.5}
\dcref{I.4.5}
For varieties $X,Y$, the following are equivalent: $X,Y$ are birational; they have isomorphic
dense open subsets; and $K(X)\cong K(Y)$ over $k$.
\uses{thm:ha-ch1-4-4}
\end{corollary}
\begin{proof}
Restrict inverse rational maps to common domains for the first implication; an isomorphism of opens
induces an isomorphism of function fields; Theorem 4.4 gives the converse.
\end{proof}

\begin{theorem}[Primitive element]
\label{thm:ha-ch1-4-6A}
\source{hartshorne-algebraic-geometry:I.4.6A}
\dcref{I.4.6A}
A finite separable extension $L/K$ is generated by one element; over an infinite $K$, a generator
can be chosen as a $K$-linear combination of any finite generating set.
\end{theorem}
\begin{proof}
Hartshorne cites Zariski--Samuel.
\end{proof}

\begin{theorem}[Separating transcendence bases]
\label{thm:ha-ch1-4-7A}
\source{hartshorne-algebraic-geometry:I.4.7A}
\dcref{I.4.7A}
Every finitely generated separably generated extension has, inside any generating set, a separating
transcendence basis.
\end{theorem}
\begin{proof}
Hartshorne cites Zariski--Samuel.
\end{proof}

\begin{theorem}[Perfect fields]
\label{thm:ha-ch1-4-8A}
\source{hartshorne-algebraic-geometry:I.4.8A}
\dcref{I.4.8A}
Every finitely generated extension of a perfect field is separably generated.
\end{theorem}
\begin{proof}
Hartshorne cites Zariski--Samuel and Matsumura.
\end{proof}

\begin{proposition}[Birational hypersurface model]
\label{prop:ha-ch1-4-9}
\source{hartshorne-algebraic-geometry:I.4.9}
\dcref{I.4.9}
Every $r$-dimensional variety is birational to a hypersurface in $\mathbb A^{r+1}$ (and to its
projective closure in $\mathbb P^{r+1}$).
\uses{cor:ha-ch1-4-5}
\end{proposition}
\begin{proof}
Choose a separating transcendence basis $x_1,\ldots,x_r$ and, by the primitive-element theorem,
$y$ with $K(X)=k(x_1,\ldots,x_r,y)$.  Clearing denominators in the algebraic equation of $y$
gives an irreducible equation in $r+1$ variables whose hypersurface has function field $K(X)$.
Apply Corollary 4.5.
\end{proof}

\section{Nonsingular Varieties}

\begin{definition}[Local, cotangent, and tangent spaces]
\label{def:ha-ch1-tangent-space}
\source{hartshorne-algebraic-geometry:I.5}
For a point $P$ of an affine variety, let $\mathfrak m_P$ be the maximal ideal of its local ring.
The cotangent space is $\mathfrak m_P/\mathfrak m_P^2$ and the Zariski tangent space is its dual.
\end{definition}

\begin{definition}[Regular local ring and nonsingularity]
\label{def:ha-ch1-regular-local}
\source{hartshorne-algebraic-geometry:I.5}
A noetherian local ring $(A,\mathfrak m,k)$ is regular if
$\dim_k\mathfrak m/\mathfrak m^2=\dim A$. A variety is nonsingular at $P$ when its local ring is
regular, and singular otherwise.
\end{definition}

\begin{definition}[Multiplicity]
\label{def:ha-ch1-multiplicity}
\source{hartshorne-algebraic-geometry:I.5}
The multiplicity of a local ring is the normalized leading coefficient of its Hilbert--Samuel
polynomial; at a point of a variety it is the multiplicity of its local ring.
\end{definition}

\begin{theorem}[Jacobian criterion]
\label{thm:ha-ch1-5-1}
\source{hartshorne-algebraic-geometry:I.5.1}
\dcref{I.5.1}
For an affine variety $Y\subset\mathbb A^n$, $P$ is nonsingular iff $\mathcal O_{P,Y}$ is
regular.  If $\dim Y=r$, this is equivalent to the Jacobian of generators of $I(Y)$ having rank
$n-r$ at $P$.
\end{theorem}
\begin{proof}
Translate $P=(a_i)$ to the origin.  The differential map identifies the cotangent space of affine
space with $k^n$, and the image of $I(Y)$ has dimension equal to the Jacobian rank.  Thus
\(\dim_k\mathfrak m_P/\mathfrak m_P^2+\operatorname{rank}J=n\).  Since
\(\dim\mathcal O_{P,Y}=r\), regularity is equivalent to rank $J=n-r$.
\end{proof}

\begin{proposition}[Embedding-dimension inequality]
\label{prop:ha-ch1-5-2A}
\source{hartshorne-algebraic-geometry:I.5.2A}
\dcref{I.5.2A}
For a noetherian local ring $(A,\mathfrak m,k)$, \(\dim_k\mathfrak m/\mathfrak m^2\ge\dim A\), with equality precisely when $A$ is regular.
\end{proposition}
\begin{proof}
Hartshorne cites Atiyah--Macdonald and Matsumura.
\end{proof}

\begin{theorem}[Openness of the nonsingular locus]
\label{thm:ha-ch1-5-3}
\source{hartshorne-algebraic-geometry:I.5.3}
\dcref{I.5.3}
For every variety $Y$, $\operatorname{Sing}Y$ is a proper closed subset.
\end{theorem}
\begin{proof}
On an affine chart of dimension $r$ it is cut out by $I(Y)$ and all $(n-r)\times(n-r)$ Jacobian
minors, hence is closed.  Birationality reduces properness to an irreducible hypersurface
$f=0$.  If every point were singular, all partial derivatives would vanish on $Y$, hence would
be multiples of $f$; their lower degrees force them to be zero.  In characteristic zero this is
impossible.  In characteristic $p$, $f$ is a polynomial in the $p$th powers of the variables,
and algebraic closedness makes $f$ a $p$th power, contradicting irreducibility.
\end{proof}

\begin{theorem}[Completion of local rings]
\label{thm:ha-ch1-5-4A}
\source{hartshorne-algebraic-geometry:I.5.4A}
\dcref{I.5.4A}
For a noetherian local ring $A$ with completion $\widehat A$: $\widehat A$ is local with maximal
ideal $\mathfrak m\widehat A$ and $A\hookrightarrow\widehat A$; the completion of every finite
$A$-module $M$ is $M\otimes_A\widehat A$; \(\dim\widehat A=\dim A\); and $A$ is regular iff
$\widehat A$ is regular.
\end{theorem}
\begin{proof}
Hartshorne cites Atiyah--Macdonald and Zariski--Samuel.
\end{proof}

\begin{theorem}[Cohen structure theorem]
\label{thm:ha-ch1-5-5A}
\source{hartshorne-algebraic-geometry:I.5.5A}
\dcref{I.5.5A}
If a complete regular local ring of dimension $n$ contains a field, then it is isomorphic to
$k[[x_1,\ldots,x_n]]$ over its residue field $k$.
\end{theorem}
\begin{proof}
Hartshorne cites Matsumura and Zariski--Samuel.
\end{proof}

\begin{theorem}[Elimination theory]
\label{thm:ha-ch1-5-7A}
\source{hartshorne-algebraic-geometry:I.5.7A}
\dcref{I.5.7A}
For homogeneous polynomials $f_1,\ldots,f_r$ with indeterminate coefficients, there are universal
integer-coefficient polynomials in those coefficients, homogeneous in the coefficients of each
$f_i$, whose simultaneous vanishing is equivalent over every field to the existence of a common
nonzero projective zero.
\end{theorem}
\begin{proof}
Hartshorne cites van der Waerden.
\end{proof}

\section{Nonsingular Curves}

\begin{definition}[Curves, valuation rings, and DVRs]
\label{def:ha-ch1-curve-valuation}
\source{hartshorne-algebraic-geometry:I.6}
A curve is a one-dimensional variety. A valuation ring of a field $K$ is a local subring $R$
such that for every $x\in K^*$ either $x$ or $x^{-1}$ lies in $R$. It is discrete when its value
group is $\mathbb Z$, in which case it is a discrete valuation ring (DVR).
\end{definition}

\begin{definition}[Domination and abstract nonsingular curve]
\label{def:ha-ch1-abstract-curve}
\source{hartshorne-algebraic-geometry:I.6}
A local ring $R$ dominates a local ring $A$ when $A\subseteq R$ and the maximal ideal of $R$
contracts to that of $A$. An abstract nonsingular curve is the curve obtained by gluing the local
models associated to the DVRs of its function field.
\end{definition}

\begin{theorem}[Valuation rings]
\label{thm:ha-ch1-6-1A}
\source{hartshorne-algebraic-geometry:I.6.1A}
\dcref{I.6.1A}
A local subring $R\subset K$ is a valuation ring iff for each $x\in K^*$, $x\in R$ or
$x^{-1}\in R$.
\end{theorem}
\begin{proof}
Hartshorne cites Bourbaki and Atiyah--Macdonald.
\end{proof}

\begin{theorem}[One-dimensional normal local rings]
\label{thm:ha-ch1-6-2A}
\source{hartshorne-algebraic-geometry:I.6.2A}
\dcref{I.6.2A}
A one-dimensional noetherian local domain is a DVR iff it is integrally closed.
\end{theorem}
\begin{proof}
Hartshorne cites Atiyah--Macdonald.
\end{proof}

\begin{theorem}[Integral closure of a Dedekind domain]
\label{thm:ha-ch1-6-3A}
\source{hartshorne-algebraic-geometry:I.6.3A}
\dcref{I.6.3A}
The integral closure of a Dedekind domain in a finite extension of its fraction field is Dedekind.
\end{theorem}
\begin{proof}
Hartshorne cites Zariski--Samuel.
\end{proof}

\begin{lemma}[Incomparability of curve local rings]
\label{lem:ha-ch1-6-4}
\source{hartshorne-algebraic-geometry:I.6.4}
\dcref{I.6.4}
If $P\ne Q$ are points of a quasi-projective variety with the same function field, neither local
ring dominates the other.
\end{lemma}
\begin{proof}
Embed in projective space and choose a homogeneous coordinate ratio regular at one point and
having a pole at the other.  It belongs to one local ring but not the other, ruling out domination.
\end{proof}

\begin{lemma}[Valuation-ring intersection]
\label{lem:ha-ch1-6-5}
\source{hartshorne-algebraic-geometry:I.6.5}
\dcref{I.6.5}
For a one-dimensional function field $K/k$ and $x\in K^*$, $x$ belongs to a valuation ring $R$
iff it has nonnegative valuation there; if $x\notin R$, then $x^{-1}$ lies in the maximal ideal.
The intersection of the valuation rings containing $k$ is the relative constant field.
\end{lemma}
\begin{proof}
The membership equivalence is the defining dichotomy of a valuation ring.  Applying it to $x^{-1}$
gives the maximal-ideal assertion; intersecting over all such rings gives the constant-field claim.
\end{proof}

\begin{corollary}[DVRs are curve local rings]
\label{cor:ha-ch1-6-6}
\source{hartshorne-algebraic-geometry:I.6.6}
\dcref{I.6.6}
Every DVR of a one-dimensional function field $K/k$ is isomorphic to the local ring of a point on
a nonsingular affine curve.
\uses{lem:ha-ch1-6-5}
\end{corollary}
\begin{proof}
Choose $y\in R\setminus k$.  The construction in Lemma 6.5 gives a finitely generated normal
one-dimensional $k$-algebra whose localization at a maximal ideal is $R$.
\end{proof}

\begin{proposition}[Abstract model of a nonsingular curve]
\label{prop:ha-ch1-6-7}
\source{hartshorne-algebraic-geometry:I.6.7}
\dcref{I.6.7}
Every nonsingular quasi-projective curve is isomorphic to an abstract nonsingular curve formed from
the DVRs of its function field.
\uses{lem:ha-ch1-6-4,cor:ha-ch1-6-6}
\end{proposition}
\begin{proof}
Map each point to its local DVR.  The map is injective by Lemma 6.4 and surjective onto the chosen
open DVR set by Corollary 6.6.  Affine neighborhoods and the valuation criterion identify the
topology and local rings, hence give an isomorphism.
\end{proof}

\begin{proposition}[Extension across finitely many points]
\label{prop:ha-ch1-6-8}
\source{hartshorne-algebraic-geometry:I.6.8}
\dcref{I.6.8}
A morphism from an open subset of a nonsingular curve to a nonsingular complete curve extends
uniquely across the finitely many omitted points.
\uses{lem:ha-ch1-4-1}
\end{proposition}
\begin{proof}
Embed the target in projective space and write the map by homogeneous coordinates in the function
field.  At an omitted DVR choose a coordinate of minimal valuation; after rescaling all coordinates
lie in the DVR and one is a unit, so they define a point and a morphism there.  Uniqueness follows
from Lemma 4.1.
\end{proof}

\begin{theorem}[Complete nonsingular model]
\label{thm:ha-ch1-6-9}
\source{hartshorne-algebraic-geometry:I.6.9}
\dcref{I.6.9}
Every one-dimensional function field $K/k$ has a nonsingular projective curve $C_K$ whose points
are exactly the DVRs of $K/k$; it is unique up to unique isomorphism.
\uses{lem:ha-ch1-6-4,prop:ha-ch1-6-8}
\end{theorem}
\begin{proof}
Choose finitely many affine normal curve models whose local rings cover all DVRs, and glue them
along common principal affine opens.  The result is nonsingular and, by the finite affine cover
and the projective-coordinate construction, projective.  Any other model has the same DVR points;
Lemma 6.4 identifies them uniquely and Lemma 6.8 extends the resulting birational map everywhere.
\end{proof}

\begin{corollary}[Projective completion of curves]
\label{cor:ha-ch1-6-10}
\source{hartshorne-algebraic-geometry:I.6.10}
\dcref{I.6.10}
Every abstract nonsingular curve is quasi-projective, and every nonsingular quasi-projective curve is
an open subset of a nonsingular projective curve.
\uses{thm:ha-ch1-6-9,prop:ha-ch1-6-7}
\end{corollary}
\begin{proof}
The construction in Theorem~\ref{thm:ha-ch1-6-9} realizes the abstract curve attached to its
function field as an open subset of the complete nonsingular model.  Conversely, Proposition
~\ref{prop:ha-ch1-6-7} identifies a nonsingular quasi-projective curve with the corresponding open
subset of its DVR space.  The omitted set is finite because a quasi-projective curve has a finite
affine cover and the complete model is proper.
\end{proof}

\begin{corollary}[Nonsingular projective model]
\label{cor:ha-ch1-6-11}
\source{hartshorne-algebraic-geometry:I.6.11}
\dcref{I.6.11}
Every curve is birationally equivalent to a nonsingular projective curve.
\uses{thm:ha-ch1-6-9}
\end{corollary}
\begin{proof}
Apply Theorem 6.9 to its one-dimensional function field.
\end{proof}

\begin{corollary}[Equivalent categories of curves]
\label{cor:ha-ch1-6-12}
\source{hartshorne-algebraic-geometry:I.6.12}
\dcref{I.6.12}
The categories of nonsingular projective curves with dominant morphisms, quasi-projective curves
with dominant rational maps, and one-dimensional function fields over $k$ with $k$-homomorphisms
are equivalent.
\uses{prop:ha-ch1-6-8,thm:ha-ch1-6-9}
\end{corollary}
\begin{proof}
Use the function-field functor and the construction $K\mapsto C_K$.  Proposition 6.8 extends the
rational maps to complete models, and uniqueness in Theorem 6.9 proves that the composites are
identity functors up to canonical isomorphism.
\end{proof}

\section{Intersections in Projective Space}

\begin{proposition}[Affine dimension theorem]
\label{prop:ha-ch1-7-1}
\source{hartshorne-algebraic-geometry:I.7.1}
\dcref{I.7.1}
If $Y,Z\subset\mathbb A^n$ have dimensions $r,s$, every irreducible component of $Y\cap Z$ has
dimension at least $r+s-n$.
\end{proposition}
\begin{proof}
For a hypersurface, minimal primes over its equation have height one by Hauptidealsatz, giving
dimension $r-1$.  For general $Z$, identify $Y\cap Z$ with $(Y\times Z)\cap\Delta$ in
$\mathbb A^{2n}$; the diagonal is cut out by $n$ hypersurfaces, so iterating the hypersurface case
gives $(r+s)-n$.
\end{proof}

\begin{theorem}[Projective dimension theorem]
\label{thm:ha-ch1-7-2}
\source{hartshorne-algebraic-geometry:I.7.2}
\dcref{I.7.2}
For varieties $Y,Z\subset\mathbb P^n$ of dimensions $r,s$, every component of $Y\cap Z$ has
dimension at least $r+s-n$; if $r+s-n\ge0$, then $Y\cap Z\ne\varnothing$.
\uses{prop:ha-ch1-7-1}
\end{theorem}
\begin{proof}
The dimension assertion follows on affine charts from Proposition 7.1.  For nonemptiness, take
the affine cones, of dimensions $r+1,s+1$, which meet at the origin.  Proposition 7.1 gives their
intersection dimension at least $r+s-n+1>0$, so it contains a nonzero point, whose projective class
lies in $Y\cap Z$.
\end{proof}

\begin{definition}[Numerical and Hilbert polynomials]
\label{def:ha-ch1-hilbert-polynomial}
\source{hartshorne-algebraic-geometry:I.7}
A numerical polynomial is a polynomial taking integer values at all sufficiently large integers.
For a finite graded module, its Hilbert polynomial is the numerical polynomial agreeing eventually
with the Hilbert function.
\end{definition}

\begin{proposition}[Numerical polynomials]
\label{prop:ha-ch1-7-3}
\source{hartshorne-algebraic-geometry:I.7.3}
\dcref{I.7.3}
(a) Every numerical polynomial has a unique integer expansion \(P(z)=\sum_{i=0}^r c_i\binom zi\).
(b) If a function $f:\mathbb Z\to\mathbb Z$ has eventually numerical-polynomial first difference,
then $f$ agrees eventually with a numerical polynomial.
\end{proposition}
\begin{proof}
Induct on degree.  Since \(\Delta\binom zi=\binom z{i-1}\), the coefficients below the top
degree are integral by induction; eventual integrality of $P$ forces the top coefficient integral.
For (b), expand the difference in the binomial basis, sum it to a polynomial $P$, and observe that
$f-P$ is eventually constant.
\end{proof}

\begin{proposition}[Graded-module filtration]
\label{prop:ha-ch1-7-4}
\source{hartshorne-algebraic-geometry:I.7.4}
\dcref{I.7.4}
Every finite graded module over a noetherian graded ring has a finite graded filtration whose factors
are shifts of $S/\mathfrak p_i$ for homogeneous primes.  Every prime containing $\operatorname{Ann}M$
contains some $\mathfrak p_i$; the minimal such primes are the minimal primes of $M$, and each occurs
with multiplicity equal to $\length_{S_\mathfrak p}M_\mathfrak p$.
\end{proposition}
\begin{proof}
Choose a maximal graded submodule with such a filtration.  In the quotient choose a nonzero
homogeneous element with maximal annihilator; the annihilator is prime, since $abm=0$ and $bm\ne0$
force $a$ into the maximal annihilator.  Adjoining the cyclic submodule extends the filtration,
contradicting maximality unless it is all of $M$.  Localizing at a minimal prime leaves precisely
the factors with that prime, each contributing the residue field, proving the multiplicity claim.
\end{proof}

\begin{theorem}[Hilbert--Serre]
\label{thm:ha-ch1-7-5}
\source{hartshorne-algebraic-geometry:I.7.5}
\dcref{I.7.5}
For a finite graded module $M$ over $S=k[x_0,\ldots,x_n]$, the Hilbert function agrees for all
large integers with a unique polynomial $P_M(z)$, and \(\deg P_M=\dim Z(\operatorname{Ann}M)\).
\uses{prop:ha-ch1-7-4,thm:ha-ch1-7-2,prop:ha-ch1-7-3}
\end{theorem}
\begin{proof}
Use the filtration of Proposition 7.4 and induction on the dimension of the support.  For
$M=S/\mathfrak p$ choose $x_i\notin\mathfrak p$ and use
$0\to M(-1)\xrightarrow{x_i}M\to M/x_iM\to0$; the last support is the hyperplane section,
of one lower dimension by Theorem 7.2.  Proposition 7.3 integrates the first differences.  Additivity
of Hilbert functions over short exact sequences and the filtration prove the general case; eventual
polynomial equality gives uniqueness.
\end{proof}

\begin{definition}[Degree]
\label{def:ha-ch1-degree}
\source{hartshorne-algebraic-geometry:I.7}
The degree of a projective variety is the leading coefficient of its Hilbert polynomial multiplied
by the factorial of its dimension.
\end{definition}

\begin{proposition}[Degree]
\label{prop:ha-ch1-7-6}
\source{hartshorne-algebraic-geometry:I.7.6}
\dcref{I.7.6}
For nonempty $Y\subset\mathbb P^n$, $\deg Y$ is a positive integer.  If $Y=Y_1\cup Y_2$ has
same-dimensional components with smaller-dimensional intersection, degrees add; $\deg\mathbb P^n=1$;
and a hypersurface defined by a degree-$d$ form has degree $d$.
\uses{prop:ha-ch1-7-3}
\end{proposition}
\begin{proof}
The binomial expansion of Proposition 7.3 makes the leading coefficient integral and positive.
The exact sequence
\(0\to S/(I_1\cap I_2)\to S/I_1\oplus S/I_2\to S/(I_1+I_2)\to0\)
shows additivity because the last term has lower-dimensional support.  Since
$P_{\mathbb P^n}=\binom{z+n}{n}$, its degree is one.  Finally
$0\to S(-d)\xrightarrow f S\to S/(f)\to0$ gives
$P_{S/(f)}(z)=\binom{z+n}{n}-\binom{z-d+n}{n}$, whose degree is $d$.
\end{proof}

\begin{definition}[Intersection multiplicity]
\label{def:ha-ch1-intersection-multiplicity}
\source{hartshorne-algebraic-geometry:I.7}
The intersection multiplicity of a component in a hypersurface intersection is the length of the
corresponding localized graded module.
\end{definition}

\begin{theorem}[Intersection with a hypersurface]
\label{thm:ha-ch1-7-7}
\source{hartshorne-algebraic-geometry:I.7.7}
\dcref{I.7.7}
Let $Y\subset\mathbb P^n$ have dimension at least one and let $H$ be a hypersurface not containing
$Y$.  If $Z_1,\ldots,Z_s$ are the components of $Y\cap H$, and
$i(Y,H;Z_j)$ is the multiplicity of the corresponding minimal prime, then
\[
\sum_j i(Y,H;Z_j)\deg Z_j=(\deg Y)(\deg H).
\]
\uses{prop:ha-ch1-7-4}
\end{theorem}
\begin{proof}
For a degree-$d$ equation of $H$, the exact sequence
$0\to(S/I_Y)(-d)\to S/I_Y\to S/(I_Y+I_H)\to0$ gives
$P_M(z)=P_Y(z)-P_Y(z-d)$.  Comparing leading coefficients yields
$d\deg Y/(r-1)!$.  Filter $M$ as in Proposition 7.4; only minimal primes of dimension $r-1$
contribute to the leading coefficient, each with multiplicity $i(Y,H;Z_j)$ and degree $\deg Z_j$.
Equating the coefficients proves the formula.
\end{proof}

\begin{corollary}[Bezout]
\label{cor:ha-ch1-7-8}
\source{hartshorne-algebraic-geometry:I.7.8}
\dcref{I.7.8}
Distinct plane curves of degrees $d,e$ with finite intersection points $P_1,\ldots,P_s$ satisfy
$\sum_i i(Y,Z;P_i)=de$.
\uses{thm:ha-ch1-7-7}
\end{corollary}
\begin{proof}
Apply Theorem 7.7.  Each point has Hilbert polynomial $1$ and degree one.
\end{proof}

\section{What Is Algebraic Geometry?}

This final section is exposition: classification by birational equivalence and isomorphism, the
curve genus/moduli picture, intrinsic invariants, cohomology, arbitrary ground fields, abstract
varieties, reducible algebraic sets, and the transition to schemes.  It contains no numbered
definition, lemma, proposition, theorem, or corollary to include in the blueprint.
